% ==============================================================================
% SECCIÓN 6: RELACIONES CON EL CLIENTE
% ==============================================================================
\section{Relaciones con el Cliente (Entrevistas, Reuniones, Revisiones)}

\subsection{Dinámica de Colaboración Semanal y Entrevistas}
Se realizará una visita al colegio para obtener la información extra necesaria no aportada en la entrevista con el profesorado. La información adicional que necesitemos se obtendrá mediante el correo de contacto proporcionado por el cliente y las reuniones con el profesorado. El objetivo primordial será comprobar la evolución del proyecto, validar las funcionalidades desarrolladas y resolver posibles dudas.  

También se propone utilizar las sesiones semanales de prácticas como reuniones con el cliente (el profesor de prácticas de la asignatura) para revisar el avance del proyecto y validar funcionalidades.

Los cambios y decisiones acordados en estas reuniones se registrarán metódicamente para mantener actualizada la planificación y los requisitos del proyecto en todo momento.

\subsection{Interlocutores y Ética en el Trato}
\begin{itemize}[leftmargin=*]
  \item \textbf{Entorno de Aplicación:} Personal directivo, docentes y profesionales del \textbf{C.E.E. Fundación Purísima Concepción de Granada}.
  \item \textbf{Destinatarios Finales:} Alumnado del Programa de Transición a la Vida Adulta (PTVAL), tutores escolares y familias.
  \item \textbf{Protección de Datos y RGPD:} Se observará máxima cautela y cumplimiento de las leyes de protección de datos de menores y personas con discapacidad, garantizando anonimato en datos y fotografías.
\end{itemize}
